\scanpage{76}
\mspage{75}

Pour toute partie $A$ de $M$, posons
\[
|A|=\bigcup_{X\in A}|X|
\]
% slanted
\margin{\uline{NB} $|\varnothing|=\varnothing$, et si $\mathrm{vi}_M$ \tikz[baseline=(n.base)]{\node[draw,circle,inner sep=0.7pt] (n) {\unsure{a}};} \unsure{vérifié|vérifiée}, $|M|=P$}

Ceci posé, je \deleted{dis} \inserted{voudrais} que pour deux parties $A,B$ de $M$, on ait
\[
\begin{aligned}
\operatorname{supp}A\subset\operatorname{supp}B
&\overset{?}{\Longleftrightarrow}|A|\subset|B|\\
\text{(ou encore }A\subset\operatorname{supp}B\text{)}&
\end{aligned}
\tag*{(*)}
\]
En effet, \deleted{\unsure{on}} \unsure{la première|le premier} \deleted{\unsure{relation}} \inserted{membre} équivaut à
\[
\operatorname{cosupp}B\subset\operatorname{cosupp}A,
\tag*{\unsure{(a)}}
\]
or \deleted{on a} on a, pour $X\in M$
\[
\begin{aligned}
X\in\operatorname{cosupp}A
&\overset{\mathrm{def}}{\Longleftrightarrow}\deleted{\unsure{A}}X\parallel A
\overset{\mathrm{def}}{\Longleftrightarrow}\forall Y\in A,\ X\parallel Y\\
&\overset{\mathrm{iii}_M}{\Longleftrightarrow}\forall Y\in A,\ |X|\cap|Y|=\varnothing
\Longleftrightarrow |X|\cap|A|=\varnothing
\end{aligned}
\]

Donc l'inclusion $\operatorname{supp}A\subset\operatorname{supp}B$, ou encore \unsure{(a)|(*)}, équivaut à
\[
\forall X\in M,\quad |X|\cap|B|=\varnothing\quad\Longrightarrow\quad |X|\cap|A|=\varnothing.
\tag*{($\bullet$)}
\]
et cette \deleted{\unsure{relation}} \inserted{propriété} sur $(A,B)$ est évidemment impliquée par $|A|\subset|B|$. Donc dans (*) on a $\Longleftarrow$. Pour voir $\Longrightarrow$, il revient au $=$ de montrer que \deleted{\unsure{les}} $|A|\not\subset|B|$ implique l'existence de $X\in M$, avec $|X|\cap|B|=\varnothing$ mais $|X|\cap|A|\ne\varnothing$. Ce n'est pas clair en général, et je \unsure{donnerai} \unsure{peut-être|peut être} contre-exemple \unsure{tantôt}. Mais supposons que $B$ soit un support, i.e.\ \deleted{$\unsure{B}$}
\[
B=\operatorname{supp}B,
\]
